\documentclass[aspectratio=169]{beamer}
\usepackage[utf8]{inputenc}
\usepackage{amsmath, amssymb}
\usepackage{graphicx}
\usepackage{booktabs}
\usepackage{tcolorbox}

% Theme styling based on UNIB colors
\definecolor{UNIBBlue}{RGB}{0, 51, 102}
\definecolor{UNIBYellow}{RGB}{255, 204, 0}
\setbeamercolor{palette primary}{bg=UNIBBlue,fg=white}
\setbeamercolor{palette secondary}{bg=UNIBYellow,fg=black}
\setbeamercolor{structure}{fg=UNIBBlue}
\setbeamercolor{title}{fg=UNIBBlue}
\setbeamercolor{frametitle}{bg=UNIBBlue,fg=white}

\title{Optimisasi untuk Teknik Elektro}
\subtitle{Minggu 4: Analisis Sensitivitas}
\author{Ir. Novalio Daratha S.T., M.Sc., Ph.D.}
\date{Semester Ganjil 2026}

\begin{document}

\begin{frame}
    \titlepage
\end{frame}

\begin{frame}{Tujuan Pembelajaran (CPMK-1, CPMK-4)}
    \begin{itemize}
        \item Mampu menjelaskan konsep dan tujuan dari Analisis Sensitivitas (Sensitivity Analysis/Post-Optimality Analysis).
        \item Mampu menentukan rentang optimalitas (Range of Optimality) untuk koefisien fungsi tujuan.
        \item Mampu menentukan rentang kelayakan (Range of Feasibility) untuk nilai ruas kanan kendala.
        \item Mampu menginterpretasikan laporan sensitivitas dari \textit{solver} komputasional untuk pengambilan keputusan di keteknikan.
    \end{itemize}
\end{frame}

\begin{frame}{Mengapa Analisis Sensitivitas?}
    \textbf{Dunia nyata selalu berubah.} Model program linier didasarkan pada parameter yang diestimasi atau diprediksi, yang bisa saja tidak akurat atau berubah di masa depan.
    \vspace{0.3cm}
    \begin{itemize}
        \item Bagaimana jika harga komponen pemancar radio naik? Apakah desain alokasi daya saat ini masih merupakan yang terbaik?
        \item Bagaimana jika kapasitas bandwidth backhaul diturunkan oleh provider? Seberapa drastis performa jaringan akan turun?
    \end{itemize}
    \textbf{Analisis Sensitivitas} menjawab pertanyaan "Bagaimana jika..." (What-if analysis) \textbf{tanpa} harus menyelesaikan ulang model dari awal.
\end{frame}

\begin{frame}{Dua Jenis Analisis Utama}
    \begin{columns}
        \begin{column}{0.5\textwidth}
            \begin{tcolorbox}[colback=blue!5,colframe=UNIBBlue,title=Perubahan Koefisien Tujuan ($c_j$)]
                \textbf{Rentang Optimalitas} \\
                Mencari seberapa besar koefisien profit/biaya suatu variabel dapat berubah sebelum solusi basis (titik sudut optimal) \textit{berubah}. \\
                (Nilai $Z$ akan berubah, tetapi strategi $X$ yang dipilih tetap sama).
            \end{tcolorbox}
        \end{column}
        \begin{column}{0.5\textwidth}
            \begin{tcolorbox}[colback=yellow!5,colframe=UNIBYellow,title=Perubahan Sisi Kanan ($b_i$)]
                \textbf{Rentang Kelayakan} \\
                Mencari seberapa besar kapasitas/batasan kendala dapat diubah dengan mempertahankan Shadow Price yang konstan. \\
                (Titik optimal \textit{pasti berubah}, tetapi kombinasi variabel aktif dalam basis tetap sama).
            \end{tcolorbox}
        \end{column}
    \end{columns}
\end{frame}

\begin{frame}{Membaca Laporan Sensitivitas (Contoh)}
    \begin{table}[]
    \centering
    \begin{tabular}{@{}llccc@{}}
    \toprule
    \textbf{Variabel} & \textbf{Final Value} & \textbf{Obj. Coef.} & \textbf{Allowable Increase} & \textbf{Allowable Decrease} \\ \midrule
    $x_1$ (Daya Gen 1) & 100 MW & \$10 & 2.5 & 4.0 \\
    $x_2$ (Daya Gen 2) & 50 MW & \$15 & 1E+30 ($\infty$) & 2.5 \\ \bottomrule
    \end{tabular}
    \end{table}
    \vspace{0.2cm}
    \textbf{Interpretasi Rentang Optimalitas}: \\
    Strategi alokasi daya (Gen 1 = 100 MW, Gen 2 = 50 MW) \textbf{akan tetap optimal} selama biaya operasional Gen 1 berada di rentang: \\
    $\$10 - \$4.0 \le c_1 \le \$10 + \$2.5$ \\
    $\mathbf{\$6.0 \le c_1 \le \$12.5}$
\end{frame}

\begin{frame}{Reduced Cost}
    \textbf{Reduced Cost} berlaku untuk variabel non-basis (variabel yang nilainya 0 pada solusi optimal).
    \vspace{0.3cm}
    \begin{itemize}
        \item Ini mengukur \textbf{seberapa banyak koefisien tujuan harus diperbaiki} (misal biaya diturunkan atau profit dinaikkan) agar variabel tersebut mulai masuk ke dalam solusi optimal ($> 0$).
        \item Jika sebuah generator tidak dihidupkan (daya = 0) karena terlalu mahal, \textit{reduced cost} memberitahu kita: seberapa jauh kita harus memotong biaya operasionalnya agar layak dinyalakan.
    \end{itemize}
\end{frame}

\begin{frame}{Aplikasi di Teknik Elektro}
    \begin{itemize}
        \item \textbf{Perencanaan Investasi:} Evaluasi batasan jaringan listrik mana yang paling sensitif dan memberikan Return on Investment (ROI) tertinggi dari \textit{shadow price}-nya, asalkan penambahan kapasitas masih dalam \textit{allowable increase}.
        \item \textbf{Penjadwalan Beban (Load Dispatch):} Mengetahui ketahanan strategi alokasi daya saat ini jika terjadi fluktuasi mendadak pada harga bahan bakar (analisis koefisien objektif).
    \end{itemize}
\end{frame}

\end{document}
